\begin{frame}{Sumário}
    \tableofcontents
\end{frame}

\section{Estrutura do documento}

\begin{frame}[fragile]{Documento mínimo}
\begin{lstlisting}
\documentclass[aspectratio=169]{beamer}
\usepackage{panqueques}

\begin{document}
  \begin{frame}{Meu primeiro slide}
    Conteúdo do slide.
  \end{frame} % encerra o slide
\end{document}
\end{lstlisting}

\begin{itemize}
    \item \verb|\documentclass| escolhe a classe e as opções do documento.
    \item \verb|\usepackage| carrega recursos adicionais.
    \item \verb|\begin| e \verb|\end| delimitam ambientes.
\end{itemize}
\end{frame}

\begin{frame}[fragile]{Organização em arquivos}
O arquivo principal pode reunir seções mantidas separadamente:

\begin{lstlisting}
\begin{document}
  \input{sections/0-custom}
  \input{sections/1-introducao}
  \input{sections/3-objetivos}
\end{document}
\end{lstlisting}

\verb|\input{caminho}| insere o conteúdo de outro arquivo no ponto indicado. A extensão \texttt{.tex} pode ser omitida.
\end{frame}

\begin{frame}[fragile]{Título e dados da apresentação}
\begin{columns}[T]
\begin{column}{0.56\textwidth}
\begin{lstlisting}
\title[Titulo curto]{Titulo completo}
\author{Nome do autor}
\institute{Instituicao}
\date{\today}

\frame{\titlepage}
\end{lstlisting}
\end{column}
\begin{column}{0.41\textwidth}
\begin{itemize}
    \item O texto entre colchetes é a versão curta usada no rodapé.
    \item \verb|\today| imprime a data da compilação.
    \item \verb|\titlepage| monta a capa com os metadados.
\end{itemize}
\end{column}
\end{columns}
\end{frame}

\begin{frame}[fragile]{Seções, slides e sumário}
\begin{lstlisting}
\section{Introducao}

\begin{frame}{Contexto}
  Texto do slide.
\end{frame} % encerra o slide

\begin{frame}{Sumario}
  \tableofcontents
\end{frame} % encerra o slide
\end{lstlisting}

\verb|\section| organiza a navegação. O argumento de \verb|frame| ou o comando \verb|\frametitle| define o título do slide.
\end{frame}
